Autonomous materials discovery promises to replace blind trial-and-error with structured search, but progress depends on whether learned models can allocate scarce validation budget better than random selection. We examine this question in a reproducible offline benchmark using real \mpdata energy-above-hull labels as an oracle. Our loop starts from 400 random observations, trains a leakage-controlled \mlp surrogate on composition and coarse cell descriptors, selects 250 additional candidates per round, and repeats for eight rounds under a fixed 2,400-candidate validation budget. We compare random search, uncertainty sampling, greedy exploitation, upper-confidence-bound selection, and diversity-aware upper-confidence-bound selection across five matched seeds. The best policy, \diverseucb, found 1,933.2 stable or near-hull candidates on average, compared with 1,364.8 for random search, a gain of 568.4 discoveries with a 95\% confidence interval of [524.0, 612.8] and Holm-adjusted $p=1.05 \times 10^{-5}$. Greedy and standard \ucb also strongly outperformed random, while pure uncertainty sampling found fewer stable candidates than random. A composition-only transfer check on \gnome R2SCAN summaries improved top-100 on-hull enrichment from 0.8735 to 0.970. These results support a narrow operational claim: surrogate-guided atomic triage can reduce wasted validation effort, but it does not by itself establish new physical materials without higher-fidelity calculation, novelty audit, and synthesis evidence.
